% Nouveau EGA V (cote 25): archive cote 25, site page 139 (PDF page 140).
% typescript. Draft compared: none.
% AI-transcribed by Claude Opus 5.5 (claude-opus-5-5, adaptive thinking, effort high), Message Batches API, batch ega25b (msgbatch_01A7n5YuXubUPNK1zsepi8RG), 2026-11-13; not yet checked by eye.
\scanpage{25/139}
% NOTE: the whole typed text of this page is an extremely faint impression (very pale carbon); almost every reading below is conjectural and must be checked on the original
\unsure{Nouveau} \uline{\unsure{EGA V}}

\unsure{IV 16.}

\unsure{Un nouveau n° :} \uline{\unsure{Ensemble singulier et composantes singulières d'une fonction,}}

\uline{\unsure{et critères différentiels}}.

\unsure{Ce n° sera utilisé au par. 20 des sections hyperplanes, mais sa}

\unsure{place naturelle me semble bien au par. 16.}

\uline{\unsure{Définition}} \unsure{1.} \unsure{Soient X un préschéma régulier, f une section de} $\unsure{\mathcal{O}_X}$ \unsure{. Un}

\unsure{point} $\unsure{x\in X}$ \unsure{est dit} \uline{\unsure{zéro}} \illegible{3} \uline{\unsure{singulier}} \unsure{pour f}

\unsure{si on a} $\unsure{f_x \in \mathfrak{m}_x^2}$ \unsure{; zéro} \uline{\unsure{singulier ordinaire}} \unsure{si c'est un zéro singulier}

\unsure{et si de plus l'élément de} $\unsure{\mathfrak{m}_x^2/\mathfrak{m}_x^3 \subset \mathrm{Sym}(\mathfrak{m}_x/\mathfrak{m}_x^2)}$ \unsure{(bien défini),}

\unsure{interprété comme une forme quadratique sur le dual} $\unsure{t_x}$ \unsure{de} $\unsure{\mathfrak{m}_x/\mathfrak{m}_x^2}$ \unsure{sur} $\unsure{k(x)}$,

\unsure{est une forme non dégénérée.} \unsure{(Un zéro singulier qui n'est pas} \illegible{2}

\unsure{lier est parfois appelé} \uline{\unsure{zéro singulier ordinaire}}\unsure{).}

\uline{\unsure{Remarques}} \unsure{2.} \unsure{Si} $\unsure{x\in V(f)}$ \unsure{, alors x est un} \inserted{\unsure{non}} \unsure{zéro singulier de f ssi}

$\unsure{f_x \neq 0}$ \unsure{et x est un point} \illegible{1} \unsure{non singulier (= régulier) de} $\unsure{V(f)}$

\unsure{i.e. que x est un point régulier de} $\unsure{V(f)}$ \unsure{et} $\unsure{V(f)\neq X}$ \unsure{au voisinage de}

% NOTE: the following five lines are enclosed in a drawn rectangle crossed by faint diagonal strokes (possibly a deletion); the frame and strokes are too faint to be certain
\unsure{x. On peut prouver que si x est singulier pour f, il est}

\unsure{ordinaire ssi} $\unsure{f_x\neq 0}$ \unsure{et la multiplicité de} $\unsure{\mathcal{O}_{V(f),x}}$ \unsure{est égale à 2}

\unsure{dans ce cas on dit aussi parfois que x est un} \uline{\unsure{point singulier ordinaire}}

\unsure{du préschéma} $\unsure{V(f)}$ \unsure{(i.e.} \unsure{ssi} $\unsure{x\in V(f)}$ \unsure{est non} \illegible{3}

$\unsure{f_x\neq 0}$ \unsure{et la multiplicité de} $\unsure{\mathcal{O}_{V(f),x}}$ \unsure{est} \unsure{2)}

\uline{\unsure{Définition}} \unsure{3.} \unsure{Soient X un préschéma lisse sur S, f une}

\unsure{section de} $\unsure{\mathcal{O}_X}$ \unsure{,} $\unsure{x\in V(f)}$ \unsure{. On dit que x est un zéro singulier sur}

\uline{\unsure{singulier}} \unsure{sur S} \unsure{(resp.} \uline{\unsure{génériquement singulier ordinaire}}\unsure{) de f sur S si}

\unsure{pour toute extension} \illegible{1} \unsure{de} $\unsure{k(s)}$ \unsure{et tout point} $\unsure{x'}$ \unsure{de} $\unsure{X'}$ \unsure{au-dessus de x}

\unsure{la localisée de} $\unsure{x'}$ \unsure{du point correspondant} \unsure{et de} $\unsure{X_s}$ \unsure{est un zéro}

\unsure{singulier (resp. singulier ordinaire) de} $\unsure{f_{s}}$.

\mspage{\hand{139}}
